\begin{table}[tp]
\centering
\caption{Quadratic finite sums: comparison for \href{https://github.com/k3nfalt/longtails/blob/main/LongTails/Setups/QuadraticFiniteSum.lean}{QuadraticFiniteSum}.}\label{tab:QuadraticFiniteSum}
\parbox{\linewidth}{\footnotesize On a real Hilbert space with its Borel measurable structure, let $F(x)=n^{-1}\sum_{i=1}^n(\langle x,A_i x\rangle/2+\langle b_i,x\rangle)$, where $n\ge1$, $A_i=A_i^*$, $\|A_i\|\le L$, and $n^{-1}\sum_i A_i\succeq\mu I$, with $0<\mu\le L$. Components may be indefinite. A minimizer $x_*$ satisfies $\|x_0-x_*\|\le R$ and $\max_i\|A_ix_*+b_i\|\le G$, with $R,G\ge0$. The exact oracle returns $A_ix+b_i$; $n,L,\mu,R,G,\varepsilon$ are supplied. Target: $\mathbb E[F(x_{\rm out})-F(x_*)]\le\varepsilon$. Row \textbf{(1)} constructs constant-step random reshuffling with independent uniform permutations per epoch; the setup imposes no algorithm restriction. Comparisons below use $0<\varepsilon<LR^2$.}\par\smallskip
\begin{threeparttable}
\footnotesize
\setlength{\tabcolsep}{4pt}\renewcommand{\arraystretch}{1.25}
\begin{tabularx}{\linewidth}{@{}>{\raggedright\arraybackslash}p{0.20\linewidth}>{\raggedright\arraybackslash}p{0.32\linewidth}>{\raggedright\arraybackslash}p{0.25\linewidth}>{\raggedright\arraybackslash}X@{}}
\toprule
\textbf{Result} & \textbf{Oracle complexity} & \textbf{Novelty} & \textbf{Limitations} \\
\midrule
\textbf{(1)} \href{https://github.com/k3nfalt/longtails/blob/main/LongTails/Theorems/RandomReshufflingQuadratic.lean}{Random\allowbreak Reshuffling\allowbreak Quadratic} \textbf{(Ours)} \href{https://github.com/k3nfalt/longtails/blob/main/scripts/reviews/RandomReshufflingQuadratic.md}{(review)}
& $\displaystyle\mathcal{O}\!\left(n+\left[\frac{nL}{\mu}+\frac{L^{3/2}G}{\mu^2\sqrt\varepsilon}\right.\right.$\newline
$\displaystyle\left.\left.{}+\left(\frac{n^2L^3G^2}{\mu^4\varepsilon}\right)^{1/3}\right]\ell\right)$
& & \\
\midrule[1.2pt]
\textbf{(2)} AGD \citep[Thm.~3.18]{qfsBubeck2015}
& $\mathcal{O}(n+n\sqrt{L/\mu}\,\ell)$
& \textbf{(1)}: No established improvement.
& \textbf{(1)}: $nL/\mu$ instead of $n\sqrt{L/\mu}$, plus heterogeneity terms; needs $G$ additionally. Expected rather than deterministic accuracy. \\
\midrule
\textbf{(3)} Katyusha X \citep[Cor.~4.4, v1]{qfsAllenZhu2018v1}; finite-dimensional restriction, $b=1$, $\ell_1=\ell_2=L$
& $\mathcal{O}((n+n^{3/4}\sqrt{L/\mu})\ell)$ expected calls
& \textbf{(1)}: Deterministic budget instead of expected cost. Scope: Hilbert-space statement; the dimension extension is formal.
& \textbf{(1)}: On finite-dimensional overlap, $nL/\mu\ge n^{3/4}\sqrt{L/\mu}$, with additional $G$-dependent terms. \\
\midrule
\textbf{(4)} RandomShuffle \citep[App.~B, Lemmas~6--8, v2]{qfsAhnYunSra2020v2}; derived bound, finite-dimensional restriction
& $\displaystyle\mathcal{O}\!\left(n+\left[\frac{nL}{\mu}\max\!\left\{1,\sqrt{\frac{L}{n\mu}}\right\}\right.\right.$\newline
$\displaystyle{}+\frac{L^{3/2}G}{\mu^2\sqrt\varepsilon}$\newline
$\displaystyle\left.\left.{}+\left(\frac{n^2L^3G^2\log n}{\mu^4\varepsilon}\right)^{1/3}\right]\ell\right)$
& \textbf{(1)}: Method-specific transient gain $\sqrt{L/(n\mu)}$ when $L>36n\mu$, $G\le n\mu\sqrt{\varepsilon/L}$. Removes $(\log n)^{1/3}$ from the cubic-root term for growing $n$; no class-wide acceleration.
& \\
\bottomrule
\end{tabularx}
\begin{tablenotes}[flushleft]\footnotesize
\item[] \textbf{Abbr:} $\ell=\log(1+LR^2/\varepsilon)$; $L$ bounds component operator norms, $\mu$ is average strong convexity, $R$ is initial distance, and $G$ bounds component gradients at $x_*$. Costs count component-gradient calls, including full-gradient snapshots; one full gradient costs $n$ calls. Expectation in \textbf{(1)} is over the epoch permutations. AGD denotes accelerated gradient descent.
\end{tablenotes}
\end{threeparttable}
\end{table}
\clearpage
